% =============================================================================
% cv.tex — DOCUMENTO MAESTRO
% =============================================================================
% Plantilla basada en:
% https://github.com/darwiin/awesome-neue-latex-cv
%
% Autor: Edison Achalma (achalmaedison)
% Licencia de plantilla: CC BY-SA 4.0
%
% Este archivo NO se edita para cambiar de perfil profesional. El perfil se
% selecciona con la macro \CVprofile, definida antes de cargar este archivo:
%
%   lualatex --jobname=cv-docencia "\def\CVprofile{docencia}% =============================================================================
% cv.tex — DOCUMENTO MAESTRO
% =============================================================================
% Plantilla basada en:
% https://github.com/darwiin/awesome-neue-latex-cv
%
% Autor: Edison Achalma (achalmaedison)
% Licencia de plantilla: CC BY-SA 4.0
%
% Este archivo NO se edita para cambiar de perfil profesional. El perfil se
% selecciona con la macro \CVprofile, definida antes de cargar este archivo:
%
%   lualatex --jobname=cv-docencia "\def\CVprofile{docencia}% =============================================================================
% cv.tex — DOCUMENTO MAESTRO
% =============================================================================
% Plantilla basada en:
% https://github.com/darwiin/awesome-neue-latex-cv
%
% Autor: Edison Achalma (achalmaedison)
% Licencia de plantilla: CC BY-SA 4.0
%
% Este archivo NO se edita para cambiar de perfil profesional. El perfil se
% selecciona con la macro \CVprofile, definida antes de cargar este archivo:
%
%   lualatex --jobname=cv-docencia "\def\CVprofile{docencia}% =============================================================================
% cv.tex — DOCUMENTO MAESTRO
% =============================================================================
% Plantilla basada en:
% https://github.com/darwiin/awesome-neue-latex-cv
%
% Autor: Edison Achalma (achalmaedison)
% Licencia de plantilla: CC BY-SA 4.0
%
% Este archivo NO se edita para cambiar de perfil profesional. El perfil se
% selecciona con la macro \CVprofile, definida antes de cargar este archivo:
%
%   lualatex --jobname=cv-docencia "\def\CVprofile{docencia}\input{cv.tex}"
%
% o, desde la raíz del repositorio, con el Makefile (forma recomendada):
%
%   make docencia          # un perfil
%   make all               # los 6 perfiles → build/cv-<perfil>.pdf
%
% Perfiles disponibles (ver profiles/):
%   sector-publico | docencia | data-science | sector-financiero
%   economia       | consultoria
%
% VERSIÓN SIN ANEXOS (solo descriptivo): definiendo \CVsinAnexos se omiten
% las constancias/certificados adjuntos al final del CV:
%
%   make docencia ANEXOS=0     → build/cv-docencia-sin-anexos.pdf
%   lualatex --jobname=cv-docencia-sin-anexos
%     "\def\CVsinAnexos{1}\def\CVprofile{docencia}\input{cv.tex}"
%
% Si \CVprofile no está definida se usa el perfil por defecto:
\ifdefined\CVprofile\else\def\CVprofile{sector-publico}\fi

% Debe compilarse con LuaLaTeX (la clase usa fontspec/fuentes OpenType)
\documentclass[localFont,alternative,10pt]{yaac-another-awesome-cv}

% -----------------------------------------------------------------------------
% Configuración compartida por todos los perfiles
% -----------------------------------------------------------------------------
\input{config/personal}   % datos personales, contacto, foto, bibliografía
\input{config/settings}   % ajustes de diseño

% Carga condicional de la sección de anexos: los perfiles la invocan con
% \inputAnexos{...}; no hace nada si \CVsinAnexos está definida (ANEXOS=0).
\newcommand{\inputAnexos}[1]{\ifdefined\CVsinAnexos\else\input{#1}\fi}

\begin{document}

	% Encabezado y pie de página
	\makecvheader
	\makecvfooter
	{\textsc{\today}}                        % Fecha (izquierda)
	{\textsc{\authorFullName{} --- CV}}      % Nombre (centro)
	{\thepage}                               % Número de página (derecha)

	% Cuerpo del CV: el perfil activo decide qué secciones se incluyen
	\input{profiles/\CVprofile}

\end{document}
"
%
% o, desde la raíz del repositorio, con el Makefile (forma recomendada):
%
%   make docencia          # un perfil
%   make all               # los 6 perfiles → build/cv-<perfil>.pdf
%
% Perfiles disponibles (ver profiles/):
%   sector-publico | docencia | data-science | sector-financiero
%   economia       | consultoria
%
% VERSIÓN SIN ANEXOS (solo descriptivo): definiendo \CVsinAnexos se omiten
% las constancias/certificados adjuntos al final del CV:
%
%   make docencia ANEXOS=0     → build/cv-docencia-sin-anexos.pdf
%   lualatex --jobname=cv-docencia-sin-anexos
%     "\def\CVsinAnexos{1}\def\CVprofile{docencia}% =============================================================================
% cv.tex — DOCUMENTO MAESTRO
% =============================================================================
% Plantilla basada en:
% https://github.com/darwiin/awesome-neue-latex-cv
%
% Autor: Edison Achalma (achalmaedison)
% Licencia de plantilla: CC BY-SA 4.0
%
% Este archivo NO se edita para cambiar de perfil profesional. El perfil se
% selecciona con la macro \CVprofile, definida antes de cargar este archivo:
%
%   lualatex --jobname=cv-docencia "\def\CVprofile{docencia}\input{cv.tex}"
%
% o, desde la raíz del repositorio, con el Makefile (forma recomendada):
%
%   make docencia          # un perfil
%   make all               # los 6 perfiles → build/cv-<perfil>.pdf
%
% Perfiles disponibles (ver profiles/):
%   sector-publico | docencia | data-science | sector-financiero
%   economia       | consultoria
%
% VERSIÓN SIN ANEXOS (solo descriptivo): definiendo \CVsinAnexos se omiten
% las constancias/certificados adjuntos al final del CV:
%
%   make docencia ANEXOS=0     → build/cv-docencia-sin-anexos.pdf
%   lualatex --jobname=cv-docencia-sin-anexos
%     "\def\CVsinAnexos{1}\def\CVprofile{docencia}\input{cv.tex}"
%
% Si \CVprofile no está definida se usa el perfil por defecto:
\ifdefined\CVprofile\else\def\CVprofile{sector-publico}\fi

% Debe compilarse con LuaLaTeX (la clase usa fontspec/fuentes OpenType)
\documentclass[localFont,alternative,10pt]{yaac-another-awesome-cv}

% -----------------------------------------------------------------------------
% Configuración compartida por todos los perfiles
% -----------------------------------------------------------------------------
\input{config/personal}   % datos personales, contacto, foto, bibliografía
\input{config/settings}   % ajustes de diseño

% Carga condicional de la sección de anexos: los perfiles la invocan con
% \inputAnexos{...}; no hace nada si \CVsinAnexos está definida (ANEXOS=0).
\newcommand{\inputAnexos}[1]{\ifdefined\CVsinAnexos\else\input{#1}\fi}

\begin{document}

	% Encabezado y pie de página
	\makecvheader
	\makecvfooter
	{\textsc{\today}}                        % Fecha (izquierda)
	{\textsc{\authorFullName{} --- CV}}      % Nombre (centro)
	{\thepage}                               % Número de página (derecha)

	% Cuerpo del CV: el perfil activo decide qué secciones se incluyen
	\input{profiles/\CVprofile}

\end{document}
"
%
% Si \CVprofile no está definida se usa el perfil por defecto:
\ifdefined\CVprofile\else\def\CVprofile{sector-publico}\fi

% Debe compilarse con LuaLaTeX (la clase usa fontspec/fuentes OpenType)
\documentclass[localFont,alternative,10pt]{yaac-another-awesome-cv}

% -----------------------------------------------------------------------------
% Configuración compartida por todos los perfiles
% -----------------------------------------------------------------------------
% =============================================================================
% CONFIGURACIÓN: Información personal y de contacto
% Archivo: config/personal.tex
% Descripción: Fuente única de los datos personales del autor. Todos los
%              perfiles (profiles/*.tex) comparten esta configuración;
%              un cambio aquí se propaga a todas las variantes del CV.
% =============================================================================

% Nombre completo (usado en el pie de página y disponible para otras secciones)
\newcommand{\authorFullName}{Elmer Edison Achalma Mendoza}

% Definir el nombre del autor (Obligatorio)
% Uso: \name{<nombre>}{<apellido>}
\name{Elmer Edison}{Achalma Mendoza}

% Definir el lema del autor (Obligatorio)
% Uso: \tagline{<lema>}
\tagline{Economista | Docente | Análisis de Datos | Sector Público}

% Definir foto del autor (Opcional)
% Uso: \photo[<forma: circular, square, roundedsquare, squircle>]{<diámetro>}{<archivo>}
% La forma por defecto es circular
\photo{2.5cm}{../assets/images/profile/profile2}

% -----------------------------------------------------------------------------
% Redes sociales e información de contacto
% -----------------------------------------------------------------------------
\socialinfo{
	% LinkedIn (Opcional)
	\linkedin{achalmaedison}

	% GitHub (Opcional)
	\github{achalmed}\\

	% ORCID (Opcional) - Identificador único de investigador
	%\orcid{0000-0000-0000-0000}

	% Google Scholar (Opcional)
	%\googlescholar{<ID-google-scholar>}{Tu Nombre}

	% Sitio web personal (Opcional)
	\personalLink{achalmaedison.netlify.app}

	% Teléfono (Opcional)
	\smartphone{+51 934 179 301}\\

	% Email (Opcional)
	\email{elmer.achalma.09@unsch.edu.pe}

	% Email alternativo (Opcional)
	%\email{achalmaedison@gmail.com}\\

	% Dirección (Opcional)
	%\address{<dirección completa>}

	% Ubicación (Opcional)
	\infos{Ayacucho - Perú}

	% Redes sociales adicionales disponibles:
	% \viadeo{<usuario-viadeo>}
	% \stackoverflow{<id-stackoverflow>}
	% \stackexchange{<id-stackexchange>}
	% \twitter{<usuario-twitter>}
	% \skype{<usuario-skype>}
	% \researchgate{<usuario-researchgate>}
	% \medium{<usuario-medium>}
	% \kaggle{<usuario-kaggle>}
}

% -----------------------------------------------------------------------------
% Bibliografía (solo se imprime si el perfil activa la sección publicaciones)
% -----------------------------------------------------------------------------
\addbibresource{../bibliography/my_publications.bib}
   % datos personales, contacto, foto, bibliografía
% =============================================================================
% CONFIGURACIÓN: Ajustes de diseño y presentación
% Archivo: config/settings.tex
% Descripción: Ajustes visuales compartidos por todos los perfiles.
%              Las opciones de clase (localFont, alternative, tamaño de
%              fuente, compact) se definen en el \documentclass de cv.tex:
%
%   [localFont]      - Usar las fuentes incluidas en la carpeta fonts/
%   [alternative]    - Usar el diseño alternativo del encabezado
%   [10pt|11pt|12pt] - Tamaño de fuente (10pt por defecto)
%   [compact]        - Reducir el espacio vertical entre experiencias
% =============================================================================

% Cambiar el tamaño de la columna izquierda (por defecto 2.5cm)
% Uso: \setleftcolumnlength{<tamaño>}  — recomendado entre 1.5cm y 2.5cm
%\setleftcolumnlength{2.0cm}
   % ajustes de diseño

% Carga condicional de la sección de anexos: los perfiles la invocan con
% \inputAnexos{...}; no hace nada si \CVsinAnexos está definida (ANEXOS=0).
\newcommand{\inputAnexos}[1]{\ifdefined\CVsinAnexos\else\input{#1}\fi}

\begin{document}

	% Encabezado y pie de página
	\makecvheader
	\makecvfooter
	{\textsc{\today}}                        % Fecha (izquierda)
	{\textsc{\authorFullName{} --- CV}}      % Nombre (centro)
	{\thepage}                               % Número de página (derecha)

	% Cuerpo del CV: el perfil activo decide qué secciones se incluyen
	\input{profiles/\CVprofile}

\end{document}
"
%
% o, desde la raíz del repositorio, con el Makefile (forma recomendada):
%
%   make docencia          # un perfil
%   make all               # los 6 perfiles → build/cv-<perfil>.pdf
%
% Perfiles disponibles (ver profiles/):
%   sector-publico | docencia | data-science | sector-financiero
%   economia       | consultoria
%
% VERSIÓN SIN ANEXOS (solo descriptivo): definiendo \CVsinAnexos se omiten
% las constancias/certificados adjuntos al final del CV:
%
%   make docencia ANEXOS=0     → build/cv-docencia-sin-anexos.pdf
%   lualatex --jobname=cv-docencia-sin-anexos
%     "\def\CVsinAnexos{1}\def\CVprofile{docencia}% =============================================================================
% cv.tex — DOCUMENTO MAESTRO
% =============================================================================
% Plantilla basada en:
% https://github.com/darwiin/awesome-neue-latex-cv
%
% Autor: Edison Achalma (achalmaedison)
% Licencia de plantilla: CC BY-SA 4.0
%
% Este archivo NO se edita para cambiar de perfil profesional. El perfil se
% selecciona con la macro \CVprofile, definida antes de cargar este archivo:
%
%   lualatex --jobname=cv-docencia "\def\CVprofile{docencia}% =============================================================================
% cv.tex — DOCUMENTO MAESTRO
% =============================================================================
% Plantilla basada en:
% https://github.com/darwiin/awesome-neue-latex-cv
%
% Autor: Edison Achalma (achalmaedison)
% Licencia de plantilla: CC BY-SA 4.0
%
% Este archivo NO se edita para cambiar de perfil profesional. El perfil se
% selecciona con la macro \CVprofile, definida antes de cargar este archivo:
%
%   lualatex --jobname=cv-docencia "\def\CVprofile{docencia}\input{cv.tex}"
%
% o, desde la raíz del repositorio, con el Makefile (forma recomendada):
%
%   make docencia          # un perfil
%   make all               # los 6 perfiles → build/cv-<perfil>.pdf
%
% Perfiles disponibles (ver profiles/):
%   sector-publico | docencia | data-science | sector-financiero
%   economia       | consultoria
%
% VERSIÓN SIN ANEXOS (solo descriptivo): definiendo \CVsinAnexos se omiten
% las constancias/certificados adjuntos al final del CV:
%
%   make docencia ANEXOS=0     → build/cv-docencia-sin-anexos.pdf
%   lualatex --jobname=cv-docencia-sin-anexos
%     "\def\CVsinAnexos{1}\def\CVprofile{docencia}\input{cv.tex}"
%
% Si \CVprofile no está definida se usa el perfil por defecto:
\ifdefined\CVprofile\else\def\CVprofile{sector-publico}\fi

% Debe compilarse con LuaLaTeX (la clase usa fontspec/fuentes OpenType)
\documentclass[localFont,alternative,10pt]{yaac-another-awesome-cv}

% -----------------------------------------------------------------------------
% Configuración compartida por todos los perfiles
% -----------------------------------------------------------------------------
\input{config/personal}   % datos personales, contacto, foto, bibliografía
\input{config/settings}   % ajustes de diseño

% Carga condicional de la sección de anexos: los perfiles la invocan con
% \inputAnexos{...}; no hace nada si \CVsinAnexos está definida (ANEXOS=0).
\newcommand{\inputAnexos}[1]{\ifdefined\CVsinAnexos\else\input{#1}\fi}

\begin{document}

	% Encabezado y pie de página
	\makecvheader
	\makecvfooter
	{\textsc{\today}}                        % Fecha (izquierda)
	{\textsc{\authorFullName{} --- CV}}      % Nombre (centro)
	{\thepage}                               % Número de página (derecha)

	% Cuerpo del CV: el perfil activo decide qué secciones se incluyen
	\input{profiles/\CVprofile}

\end{document}
"
%
% o, desde la raíz del repositorio, con el Makefile (forma recomendada):
%
%   make docencia          # un perfil
%   make all               # los 6 perfiles → build/cv-<perfil>.pdf
%
% Perfiles disponibles (ver profiles/):
%   sector-publico | docencia | data-science | sector-financiero
%   economia       | consultoria
%
% VERSIÓN SIN ANEXOS (solo descriptivo): definiendo \CVsinAnexos se omiten
% las constancias/certificados adjuntos al final del CV:
%
%   make docencia ANEXOS=0     → build/cv-docencia-sin-anexos.pdf
%   lualatex --jobname=cv-docencia-sin-anexos
%     "\def\CVsinAnexos{1}\def\CVprofile{docencia}% =============================================================================
% cv.tex — DOCUMENTO MAESTRO
% =============================================================================
% Plantilla basada en:
% https://github.com/darwiin/awesome-neue-latex-cv
%
% Autor: Edison Achalma (achalmaedison)
% Licencia de plantilla: CC BY-SA 4.0
%
% Este archivo NO se edita para cambiar de perfil profesional. El perfil se
% selecciona con la macro \CVprofile, definida antes de cargar este archivo:
%
%   lualatex --jobname=cv-docencia "\def\CVprofile{docencia}\input{cv.tex}"
%
% o, desde la raíz del repositorio, con el Makefile (forma recomendada):
%
%   make docencia          # un perfil
%   make all               # los 6 perfiles → build/cv-<perfil>.pdf
%
% Perfiles disponibles (ver profiles/):
%   sector-publico | docencia | data-science | sector-financiero
%   economia       | consultoria
%
% VERSIÓN SIN ANEXOS (solo descriptivo): definiendo \CVsinAnexos se omiten
% las constancias/certificados adjuntos al final del CV:
%
%   make docencia ANEXOS=0     → build/cv-docencia-sin-anexos.pdf
%   lualatex --jobname=cv-docencia-sin-anexos
%     "\def\CVsinAnexos{1}\def\CVprofile{docencia}\input{cv.tex}"
%
% Si \CVprofile no está definida se usa el perfil por defecto:
\ifdefined\CVprofile\else\def\CVprofile{sector-publico}\fi

% Debe compilarse con LuaLaTeX (la clase usa fontspec/fuentes OpenType)
\documentclass[localFont,alternative,10pt]{yaac-another-awesome-cv}

% -----------------------------------------------------------------------------
% Configuración compartida por todos los perfiles
% -----------------------------------------------------------------------------
\input{config/personal}   % datos personales, contacto, foto, bibliografía
\input{config/settings}   % ajustes de diseño

% Carga condicional de la sección de anexos: los perfiles la invocan con
% \inputAnexos{...}; no hace nada si \CVsinAnexos está definida (ANEXOS=0).
\newcommand{\inputAnexos}[1]{\ifdefined\CVsinAnexos\else\input{#1}\fi}

\begin{document}

	% Encabezado y pie de página
	\makecvheader
	\makecvfooter
	{\textsc{\today}}                        % Fecha (izquierda)
	{\textsc{\authorFullName{} --- CV}}      % Nombre (centro)
	{\thepage}                               % Número de página (derecha)

	% Cuerpo del CV: el perfil activo decide qué secciones se incluyen
	\input{profiles/\CVprofile}

\end{document}
"
%
% Si \CVprofile no está definida se usa el perfil por defecto:
\ifdefined\CVprofile\else\def\CVprofile{sector-publico}\fi

% Debe compilarse con LuaLaTeX (la clase usa fontspec/fuentes OpenType)
\documentclass[localFont,alternative,10pt]{yaac-another-awesome-cv}

% -----------------------------------------------------------------------------
% Configuración compartida por todos los perfiles
% -----------------------------------------------------------------------------
% =============================================================================
% CONFIGURACIÓN: Información personal y de contacto
% Archivo: config/personal.tex
% Descripción: Fuente única de los datos personales del autor. Todos los
%              perfiles (profiles/*.tex) comparten esta configuración;
%              un cambio aquí se propaga a todas las variantes del CV.
% =============================================================================

% Nombre completo (usado en el pie de página y disponible para otras secciones)
\newcommand{\authorFullName}{Elmer Edison Achalma Mendoza}

% Definir el nombre del autor (Obligatorio)
% Uso: \name{<nombre>}{<apellido>}
\name{Elmer Edison}{Achalma Mendoza}

% Definir el lema del autor (Obligatorio)
% Uso: \tagline{<lema>}
\tagline{Economista | Docente | Análisis de Datos | Sector Público}

% Definir foto del autor (Opcional)
% Uso: \photo[<forma: circular, square, roundedsquare, squircle>]{<diámetro>}{<archivo>}
% La forma por defecto es circular
\photo{2.5cm}{../assets/images/profile/profile2}

% -----------------------------------------------------------------------------
% Redes sociales e información de contacto
% -----------------------------------------------------------------------------
\socialinfo{
	% LinkedIn (Opcional)
	\linkedin{achalmaedison}

	% GitHub (Opcional)
	\github{achalmed}\\

	% ORCID (Opcional) - Identificador único de investigador
	%\orcid{0000-0000-0000-0000}

	% Google Scholar (Opcional)
	%\googlescholar{<ID-google-scholar>}{Tu Nombre}

	% Sitio web personal (Opcional)
	\personalLink{achalmaedison.netlify.app}

	% Teléfono (Opcional)
	\smartphone{+51 934 179 301}\\

	% Email (Opcional)
	\email{elmer.achalma.09@unsch.edu.pe}

	% Email alternativo (Opcional)
	%\email{achalmaedison@gmail.com}\\

	% Dirección (Opcional)
	%\address{<dirección completa>}

	% Ubicación (Opcional)
	\infos{Ayacucho - Perú}

	% Redes sociales adicionales disponibles:
	% \viadeo{<usuario-viadeo>}
	% \stackoverflow{<id-stackoverflow>}
	% \stackexchange{<id-stackexchange>}
	% \twitter{<usuario-twitter>}
	% \skype{<usuario-skype>}
	% \researchgate{<usuario-researchgate>}
	% \medium{<usuario-medium>}
	% \kaggle{<usuario-kaggle>}
}

% -----------------------------------------------------------------------------
% Bibliografía (solo se imprime si el perfil activa la sección publicaciones)
% -----------------------------------------------------------------------------
\addbibresource{../bibliography/my_publications.bib}
   % datos personales, contacto, foto, bibliografía
% =============================================================================
% CONFIGURACIÓN: Ajustes de diseño y presentación
% Archivo: config/settings.tex
% Descripción: Ajustes visuales compartidos por todos los perfiles.
%              Las opciones de clase (localFont, alternative, tamaño de
%              fuente, compact) se definen en el \documentclass de cv.tex:
%
%   [localFont]      - Usar las fuentes incluidas en la carpeta fonts/
%   [alternative]    - Usar el diseño alternativo del encabezado
%   [10pt|11pt|12pt] - Tamaño de fuente (10pt por defecto)
%   [compact]        - Reducir el espacio vertical entre experiencias
% =============================================================================

% Cambiar el tamaño de la columna izquierda (por defecto 2.5cm)
% Uso: \setleftcolumnlength{<tamaño>}  — recomendado entre 1.5cm y 2.5cm
%\setleftcolumnlength{2.0cm}
   % ajustes de diseño

% Carga condicional de la sección de anexos: los perfiles la invocan con
% \inputAnexos{...}; no hace nada si \CVsinAnexos está definida (ANEXOS=0).
\newcommand{\inputAnexos}[1]{\ifdefined\CVsinAnexos\else\input{#1}\fi}

\begin{document}

	% Encabezado y pie de página
	\makecvheader
	\makecvfooter
	{\textsc{\today}}                        % Fecha (izquierda)
	{\textsc{\authorFullName{} --- CV}}      % Nombre (centro)
	{\thepage}                               % Número de página (derecha)

	% Cuerpo del CV: el perfil activo decide qué secciones se incluyen
	\input{profiles/\CVprofile}

\end{document}
"
%
% Si \CVprofile no está definida se usa el perfil por defecto:
\ifdefined\CVprofile\else\def\CVprofile{sector-publico}\fi

% Debe compilarse con LuaLaTeX (la clase usa fontspec/fuentes OpenType)
\documentclass[localFont,alternative,10pt]{yaac-another-awesome-cv}

% -----------------------------------------------------------------------------
% Configuración compartida por todos los perfiles
% -----------------------------------------------------------------------------
% =============================================================================
% CONFIGURACIÓN: Información personal y de contacto
% Archivo: config/personal.tex
% Descripción: Fuente única de los datos personales del autor. Todos los
%              perfiles (profiles/*.tex) comparten esta configuración;
%              un cambio aquí se propaga a todas las variantes del CV.
% =============================================================================

% Nombre completo (usado en el pie de página y disponible para otras secciones)
\newcommand{\authorFullName}{Elmer Edison Achalma Mendoza}

% Definir el nombre del autor (Obligatorio)
% Uso: \name{<nombre>}{<apellido>}
\name{Elmer Edison}{Achalma Mendoza}

% Definir el lema del autor (Obligatorio)
% Uso: \tagline{<lema>}
\tagline{Economista | Docente | Análisis de Datos | Sector Público}

% Definir foto del autor (Opcional)
% Uso: \photo[<forma: circular, square, roundedsquare, squircle>]{<diámetro>}{<archivo>}
% La forma por defecto es circular
\photo{2.5cm}{../assets/images/profile/profile2}

% -----------------------------------------------------------------------------
% Redes sociales e información de contacto
% -----------------------------------------------------------------------------
\socialinfo{
	% LinkedIn (Opcional)
	\linkedin{achalmaedison}

	% GitHub (Opcional)
	\github{achalmed}\\

	% ORCID (Opcional) - Identificador único de investigador
	%\orcid{0000-0000-0000-0000}

	% Google Scholar (Opcional)
	%\googlescholar{<ID-google-scholar>}{Tu Nombre}

	% Sitio web personal (Opcional)
	\personalLink{achalmaedison.netlify.app}

	% Teléfono (Opcional)
	\smartphone{+51 934 179 301}\\

	% Email (Opcional)
	\email{elmer.achalma.09@unsch.edu.pe}

	% Email alternativo (Opcional)
	%\email{achalmaedison@gmail.com}\\

	% Dirección (Opcional)
	%\address{<dirección completa>}

	% Ubicación (Opcional)
	\infos{Ayacucho - Perú}

	% Redes sociales adicionales disponibles:
	% \viadeo{<usuario-viadeo>}
	% \stackoverflow{<id-stackoverflow>}
	% \stackexchange{<id-stackexchange>}
	% \twitter{<usuario-twitter>}
	% \skype{<usuario-skype>}
	% \researchgate{<usuario-researchgate>}
	% \medium{<usuario-medium>}
	% \kaggle{<usuario-kaggle>}
}

% -----------------------------------------------------------------------------
% Bibliografía (solo se imprime si el perfil activa la sección publicaciones)
% -----------------------------------------------------------------------------
\addbibresource{../bibliography/my_publications.bib}
   % datos personales, contacto, foto, bibliografía
% =============================================================================
% CONFIGURACIÓN: Ajustes de diseño y presentación
% Archivo: config/settings.tex
% Descripción: Ajustes visuales compartidos por todos los perfiles.
%              Las opciones de clase (localFont, alternative, tamaño de
%              fuente, compact) se definen en el \documentclass de cv.tex:
%
%   [localFont]      - Usar las fuentes incluidas en la carpeta fonts/
%   [alternative]    - Usar el diseño alternativo del encabezado
%   [10pt|11pt|12pt] - Tamaño de fuente (10pt por defecto)
%   [compact]        - Reducir el espacio vertical entre experiencias
% =============================================================================

% Cambiar el tamaño de la columna izquierda (por defecto 2.5cm)
% Uso: \setleftcolumnlength{<tamaño>}  — recomendado entre 1.5cm y 2.5cm
%\setleftcolumnlength{2.0cm}
   % ajustes de diseño

% Carga condicional de la sección de anexos: los perfiles la invocan con
% \inputAnexos{...}; no hace nada si \CVsinAnexos está definida (ANEXOS=0).
\newcommand{\inputAnexos}[1]{\ifdefined\CVsinAnexos\else\input{#1}\fi}

\begin{document}

	% Encabezado y pie de página
	\makecvheader
	\makecvfooter
	{\textsc{\today}}                        % Fecha (izquierda)
	{\textsc{\authorFullName{} --- CV}}      % Nombre (centro)
	{\thepage}                               % Número de página (derecha)

	% Cuerpo del CV: el perfil activo decide qué secciones se incluyen
	\input{profiles/\CVprofile}

\end{document}
"
%
% o, desde la raíz del repositorio, con el Makefile (forma recomendada):
%
%   make docencia          # un perfil
%   make all               # los 6 perfiles → build/cv-<perfil>.pdf
%
% Perfiles disponibles (ver profiles/):
%   sector-publico | docencia | data-science | sector-financiero
%   economia       | consultoria
%
% VERSIÓN SIN ANEXOS (solo descriptivo): definiendo \CVsinAnexos se omiten
% las constancias/certificados adjuntos al final del CV:
%
%   make docencia ANEXOS=0     → build/cv-docencia-sin-anexos.pdf
%   lualatex --jobname=cv-docencia-sin-anexos
%     "\def\CVsinAnexos{1}\def\CVprofile{docencia}% =============================================================================
% cv.tex — DOCUMENTO MAESTRO
% =============================================================================
% Plantilla basada en:
% https://github.com/darwiin/awesome-neue-latex-cv
%
% Autor: Edison Achalma (achalmaedison)
% Licencia de plantilla: CC BY-SA 4.0
%
% Este archivo NO se edita para cambiar de perfil profesional. El perfil se
% selecciona con la macro \CVprofile, definida antes de cargar este archivo:
%
%   lualatex --jobname=cv-docencia "\def\CVprofile{docencia}% =============================================================================
% cv.tex — DOCUMENTO MAESTRO
% =============================================================================
% Plantilla basada en:
% https://github.com/darwiin/awesome-neue-latex-cv
%
% Autor: Edison Achalma (achalmaedison)
% Licencia de plantilla: CC BY-SA 4.0
%
% Este archivo NO se edita para cambiar de perfil profesional. El perfil se
% selecciona con la macro \CVprofile, definida antes de cargar este archivo:
%
%   lualatex --jobname=cv-docencia "\def\CVprofile{docencia}% =============================================================================
% cv.tex — DOCUMENTO MAESTRO
% =============================================================================
% Plantilla basada en:
% https://github.com/darwiin/awesome-neue-latex-cv
%
% Autor: Edison Achalma (achalmaedison)
% Licencia de plantilla: CC BY-SA 4.0
%
% Este archivo NO se edita para cambiar de perfil profesional. El perfil se
% selecciona con la macro \CVprofile, definida antes de cargar este archivo:
%
%   lualatex --jobname=cv-docencia "\def\CVprofile{docencia}\input{cv.tex}"
%
% o, desde la raíz del repositorio, con el Makefile (forma recomendada):
%
%   make docencia          # un perfil
%   make all               # los 6 perfiles → build/cv-<perfil>.pdf
%
% Perfiles disponibles (ver profiles/):
%   sector-publico | docencia | data-science | sector-financiero
%   economia       | consultoria
%
% VERSIÓN SIN ANEXOS (solo descriptivo): definiendo \CVsinAnexos se omiten
% las constancias/certificados adjuntos al final del CV:
%
%   make docencia ANEXOS=0     → build/cv-docencia-sin-anexos.pdf
%   lualatex --jobname=cv-docencia-sin-anexos
%     "\def\CVsinAnexos{1}\def\CVprofile{docencia}\input{cv.tex}"
%
% Si \CVprofile no está definida se usa el perfil por defecto:
\ifdefined\CVprofile\else\def\CVprofile{sector-publico}\fi

% Debe compilarse con LuaLaTeX (la clase usa fontspec/fuentes OpenType)
\documentclass[localFont,alternative,10pt]{yaac-another-awesome-cv}

% -----------------------------------------------------------------------------
% Configuración compartida por todos los perfiles
% -----------------------------------------------------------------------------
\input{config/personal}   % datos personales, contacto, foto, bibliografía
\input{config/settings}   % ajustes de diseño

% Carga condicional de la sección de anexos: los perfiles la invocan con
% \inputAnexos{...}; no hace nada si \CVsinAnexos está definida (ANEXOS=0).
\newcommand{\inputAnexos}[1]{\ifdefined\CVsinAnexos\else\input{#1}\fi}

\begin{document}

	% Encabezado y pie de página
	\makecvheader
	\makecvfooter
	{\textsc{\today}}                        % Fecha (izquierda)
	{\textsc{\authorFullName{} --- CV}}      % Nombre (centro)
	{\thepage}                               % Número de página (derecha)

	% Cuerpo del CV: el perfil activo decide qué secciones se incluyen
	\input{profiles/\CVprofile}

\end{document}
"
%
% o, desde la raíz del repositorio, con el Makefile (forma recomendada):
%
%   make docencia          # un perfil
%   make all               # los 6 perfiles → build/cv-<perfil>.pdf
%
% Perfiles disponibles (ver profiles/):
%   sector-publico | docencia | data-science | sector-financiero
%   economia       | consultoria
%
% VERSIÓN SIN ANEXOS (solo descriptivo): definiendo \CVsinAnexos se omiten
% las constancias/certificados adjuntos al final del CV:
%
%   make docencia ANEXOS=0     → build/cv-docencia-sin-anexos.pdf
%   lualatex --jobname=cv-docencia-sin-anexos
%     "\def\CVsinAnexos{1}\def\CVprofile{docencia}% =============================================================================
% cv.tex — DOCUMENTO MAESTRO
% =============================================================================
% Plantilla basada en:
% https://github.com/darwiin/awesome-neue-latex-cv
%
% Autor: Edison Achalma (achalmaedison)
% Licencia de plantilla: CC BY-SA 4.0
%
% Este archivo NO se edita para cambiar de perfil profesional. El perfil se
% selecciona con la macro \CVprofile, definida antes de cargar este archivo:
%
%   lualatex --jobname=cv-docencia "\def\CVprofile{docencia}\input{cv.tex}"
%
% o, desde la raíz del repositorio, con el Makefile (forma recomendada):
%
%   make docencia          # un perfil
%   make all               # los 6 perfiles → build/cv-<perfil>.pdf
%
% Perfiles disponibles (ver profiles/):
%   sector-publico | docencia | data-science | sector-financiero
%   economia       | consultoria
%
% VERSIÓN SIN ANEXOS (solo descriptivo): definiendo \CVsinAnexos se omiten
% las constancias/certificados adjuntos al final del CV:
%
%   make docencia ANEXOS=0     → build/cv-docencia-sin-anexos.pdf
%   lualatex --jobname=cv-docencia-sin-anexos
%     "\def\CVsinAnexos{1}\def\CVprofile{docencia}\input{cv.tex}"
%
% Si \CVprofile no está definida se usa el perfil por defecto:
\ifdefined\CVprofile\else\def\CVprofile{sector-publico}\fi

% Debe compilarse con LuaLaTeX (la clase usa fontspec/fuentes OpenType)
\documentclass[localFont,alternative,10pt]{yaac-another-awesome-cv}

% -----------------------------------------------------------------------------
% Configuración compartida por todos los perfiles
% -----------------------------------------------------------------------------
\input{config/personal}   % datos personales, contacto, foto, bibliografía
\input{config/settings}   % ajustes de diseño

% Carga condicional de la sección de anexos: los perfiles la invocan con
% \inputAnexos{...}; no hace nada si \CVsinAnexos está definida (ANEXOS=0).
\newcommand{\inputAnexos}[1]{\ifdefined\CVsinAnexos\else\input{#1}\fi}

\begin{document}

	% Encabezado y pie de página
	\makecvheader
	\makecvfooter
	{\textsc{\today}}                        % Fecha (izquierda)
	{\textsc{\authorFullName{} --- CV}}      % Nombre (centro)
	{\thepage}                               % Número de página (derecha)

	% Cuerpo del CV: el perfil activo decide qué secciones se incluyen
	\input{profiles/\CVprofile}

\end{document}
"
%
% Si \CVprofile no está definida se usa el perfil por defecto:
\ifdefined\CVprofile\else\def\CVprofile{sector-publico}\fi

% Debe compilarse con LuaLaTeX (la clase usa fontspec/fuentes OpenType)
\documentclass[localFont,alternative,10pt]{yaac-another-awesome-cv}

% -----------------------------------------------------------------------------
% Configuración compartida por todos los perfiles
% -----------------------------------------------------------------------------
% =============================================================================
% CONFIGURACIÓN: Información personal y de contacto
% Archivo: config/personal.tex
% Descripción: Fuente única de los datos personales del autor. Todos los
%              perfiles (profiles/*.tex) comparten esta configuración;
%              un cambio aquí se propaga a todas las variantes del CV.
% =============================================================================

% Nombre completo (usado en el pie de página y disponible para otras secciones)
\newcommand{\authorFullName}{Elmer Edison Achalma Mendoza}

% Definir el nombre del autor (Obligatorio)
% Uso: \name{<nombre>}{<apellido>}
\name{Elmer Edison}{Achalma Mendoza}

% Definir el lema del autor (Obligatorio)
% Uso: \tagline{<lema>}
\tagline{Economista | Docente | Análisis de Datos | Sector Público}

% Definir foto del autor (Opcional)
% Uso: \photo[<forma: circular, square, roundedsquare, squircle>]{<diámetro>}{<archivo>}
% La forma por defecto es circular
\photo{2.5cm}{../assets/images/profile/profile2}

% -----------------------------------------------------------------------------
% Redes sociales e información de contacto
% -----------------------------------------------------------------------------
\socialinfo{
	% LinkedIn (Opcional)
	\linkedin{achalmaedison}

	% GitHub (Opcional)
	\github{achalmed}\\

	% ORCID (Opcional) - Identificador único de investigador
	%\orcid{0000-0000-0000-0000}

	% Google Scholar (Opcional)
	%\googlescholar{<ID-google-scholar>}{Tu Nombre}

	% Sitio web personal (Opcional)
	\personalLink{achalmaedison.netlify.app}

	% Teléfono (Opcional)
	\smartphone{+51 934 179 301}\\

	% Email (Opcional)
	\email{elmer.achalma.09@unsch.edu.pe}

	% Email alternativo (Opcional)
	%\email{achalmaedison@gmail.com}\\

	% Dirección (Opcional)
	%\address{<dirección completa>}

	% Ubicación (Opcional)
	\infos{Ayacucho - Perú}

	% Redes sociales adicionales disponibles:
	% \viadeo{<usuario-viadeo>}
	% \stackoverflow{<id-stackoverflow>}
	% \stackexchange{<id-stackexchange>}
	% \twitter{<usuario-twitter>}
	% \skype{<usuario-skype>}
	% \researchgate{<usuario-researchgate>}
	% \medium{<usuario-medium>}
	% \kaggle{<usuario-kaggle>}
}

% -----------------------------------------------------------------------------
% Bibliografía (solo se imprime si el perfil activa la sección publicaciones)
% -----------------------------------------------------------------------------
\addbibresource{../bibliography/my_publications.bib}
   % datos personales, contacto, foto, bibliografía
% =============================================================================
% CONFIGURACIÓN: Ajustes de diseño y presentación
% Archivo: config/settings.tex
% Descripción: Ajustes visuales compartidos por todos los perfiles.
%              Las opciones de clase (localFont, alternative, tamaño de
%              fuente, compact) se definen en el \documentclass de cv.tex:
%
%   [localFont]      - Usar las fuentes incluidas en la carpeta fonts/
%   [alternative]    - Usar el diseño alternativo del encabezado
%   [10pt|11pt|12pt] - Tamaño de fuente (10pt por defecto)
%   [compact]        - Reducir el espacio vertical entre experiencias
% =============================================================================

% Cambiar el tamaño de la columna izquierda (por defecto 2.5cm)
% Uso: \setleftcolumnlength{<tamaño>}  — recomendado entre 1.5cm y 2.5cm
%\setleftcolumnlength{2.0cm}
   % ajustes de diseño

% Carga condicional de la sección de anexos: los perfiles la invocan con
% \inputAnexos{...}; no hace nada si \CVsinAnexos está definida (ANEXOS=0).
\newcommand{\inputAnexos}[1]{\ifdefined\CVsinAnexos\else\input{#1}\fi}

\begin{document}

	% Encabezado y pie de página
	\makecvheader
	\makecvfooter
	{\textsc{\today}}                        % Fecha (izquierda)
	{\textsc{\authorFullName{} --- CV}}      % Nombre (centro)
	{\thepage}                               % Número de página (derecha)

	% Cuerpo del CV: el perfil activo decide qué secciones se incluyen
	\input{profiles/\CVprofile}

\end{document}
"
%
% o, desde la raíz del repositorio, con el Makefile (forma recomendada):
%
%   make docencia          # un perfil
%   make all               # los 6 perfiles → build/cv-<perfil>.pdf
%
% Perfiles disponibles (ver profiles/):
%   sector-publico | docencia | data-science | sector-financiero
%   economia       | consultoria
%
% VERSIÓN SIN ANEXOS (solo descriptivo): definiendo \CVsinAnexos se omiten
% las constancias/certificados adjuntos al final del CV:
%
%   make docencia ANEXOS=0     → build/cv-docencia-sin-anexos.pdf
%   lualatex --jobname=cv-docencia-sin-anexos
%     "\def\CVsinAnexos{1}\def\CVprofile{docencia}% =============================================================================
% cv.tex — DOCUMENTO MAESTRO
% =============================================================================
% Plantilla basada en:
% https://github.com/darwiin/awesome-neue-latex-cv
%
% Autor: Edison Achalma (achalmaedison)
% Licencia de plantilla: CC BY-SA 4.0
%
% Este archivo NO se edita para cambiar de perfil profesional. El perfil se
% selecciona con la macro \CVprofile, definida antes de cargar este archivo:
%
%   lualatex --jobname=cv-docencia "\def\CVprofile{docencia}% =============================================================================
% cv.tex — DOCUMENTO MAESTRO
% =============================================================================
% Plantilla basada en:
% https://github.com/darwiin/awesome-neue-latex-cv
%
% Autor: Edison Achalma (achalmaedison)
% Licencia de plantilla: CC BY-SA 4.0
%
% Este archivo NO se edita para cambiar de perfil profesional. El perfil se
% selecciona con la macro \CVprofile, definida antes de cargar este archivo:
%
%   lualatex --jobname=cv-docencia "\def\CVprofile{docencia}\input{cv.tex}"
%
% o, desde la raíz del repositorio, con el Makefile (forma recomendada):
%
%   make docencia          # un perfil
%   make all               # los 6 perfiles → build/cv-<perfil>.pdf
%
% Perfiles disponibles (ver profiles/):
%   sector-publico | docencia | data-science | sector-financiero
%   economia       | consultoria
%
% VERSIÓN SIN ANEXOS (solo descriptivo): definiendo \CVsinAnexos se omiten
% las constancias/certificados adjuntos al final del CV:
%
%   make docencia ANEXOS=0     → build/cv-docencia-sin-anexos.pdf
%   lualatex --jobname=cv-docencia-sin-anexos
%     "\def\CVsinAnexos{1}\def\CVprofile{docencia}\input{cv.tex}"
%
% Si \CVprofile no está definida se usa el perfil por defecto:
\ifdefined\CVprofile\else\def\CVprofile{sector-publico}\fi

% Debe compilarse con LuaLaTeX (la clase usa fontspec/fuentes OpenType)
\documentclass[localFont,alternative,10pt]{yaac-another-awesome-cv}

% -----------------------------------------------------------------------------
% Configuración compartida por todos los perfiles
% -----------------------------------------------------------------------------
\input{config/personal}   % datos personales, contacto, foto, bibliografía
\input{config/settings}   % ajustes de diseño

% Carga condicional de la sección de anexos: los perfiles la invocan con
% \inputAnexos{...}; no hace nada si \CVsinAnexos está definida (ANEXOS=0).
\newcommand{\inputAnexos}[1]{\ifdefined\CVsinAnexos\else\input{#1}\fi}

\begin{document}

	% Encabezado y pie de página
	\makecvheader
	\makecvfooter
	{\textsc{\today}}                        % Fecha (izquierda)
	{\textsc{\authorFullName{} --- CV}}      % Nombre (centro)
	{\thepage}                               % Número de página (derecha)

	% Cuerpo del CV: el perfil activo decide qué secciones se incluyen
	\input{profiles/\CVprofile}

\end{document}
"
%
% o, desde la raíz del repositorio, con el Makefile (forma recomendada):
%
%   make docencia          # un perfil
%   make all               # los 6 perfiles → build/cv-<perfil>.pdf
%
% Perfiles disponibles (ver profiles/):
%   sector-publico | docencia | data-science | sector-financiero
%   economia       | consultoria
%
% VERSIÓN SIN ANEXOS (solo descriptivo): definiendo \CVsinAnexos se omiten
% las constancias/certificados adjuntos al final del CV:
%
%   make docencia ANEXOS=0     → build/cv-docencia-sin-anexos.pdf
%   lualatex --jobname=cv-docencia-sin-anexos
%     "\def\CVsinAnexos{1}\def\CVprofile{docencia}% =============================================================================
% cv.tex — DOCUMENTO MAESTRO
% =============================================================================
% Plantilla basada en:
% https://github.com/darwiin/awesome-neue-latex-cv
%
% Autor: Edison Achalma (achalmaedison)
% Licencia de plantilla: CC BY-SA 4.0
%
% Este archivo NO se edita para cambiar de perfil profesional. El perfil se
% selecciona con la macro \CVprofile, definida antes de cargar este archivo:
%
%   lualatex --jobname=cv-docencia "\def\CVprofile{docencia}\input{cv.tex}"
%
% o, desde la raíz del repositorio, con el Makefile (forma recomendada):
%
%   make docencia          # un perfil
%   make all               # los 6 perfiles → build/cv-<perfil>.pdf
%
% Perfiles disponibles (ver profiles/):
%   sector-publico | docencia | data-science | sector-financiero
%   economia       | consultoria
%
% VERSIÓN SIN ANEXOS (solo descriptivo): definiendo \CVsinAnexos se omiten
% las constancias/certificados adjuntos al final del CV:
%
%   make docencia ANEXOS=0     → build/cv-docencia-sin-anexos.pdf
%   lualatex --jobname=cv-docencia-sin-anexos
%     "\def\CVsinAnexos{1}\def\CVprofile{docencia}\input{cv.tex}"
%
% Si \CVprofile no está definida se usa el perfil por defecto:
\ifdefined\CVprofile\else\def\CVprofile{sector-publico}\fi

% Debe compilarse con LuaLaTeX (la clase usa fontspec/fuentes OpenType)
\documentclass[localFont,alternative,10pt]{yaac-another-awesome-cv}

% -----------------------------------------------------------------------------
% Configuración compartida por todos los perfiles
% -----------------------------------------------------------------------------
\input{config/personal}   % datos personales, contacto, foto, bibliografía
\input{config/settings}   % ajustes de diseño

% Carga condicional de la sección de anexos: los perfiles la invocan con
% \inputAnexos{...}; no hace nada si \CVsinAnexos está definida (ANEXOS=0).
\newcommand{\inputAnexos}[1]{\ifdefined\CVsinAnexos\else\input{#1}\fi}

\begin{document}

	% Encabezado y pie de página
	\makecvheader
	\makecvfooter
	{\textsc{\today}}                        % Fecha (izquierda)
	{\textsc{\authorFullName{} --- CV}}      % Nombre (centro)
	{\thepage}                               % Número de página (derecha)

	% Cuerpo del CV: el perfil activo decide qué secciones se incluyen
	\input{profiles/\CVprofile}

\end{document}
"
%
% Si \CVprofile no está definida se usa el perfil por defecto:
\ifdefined\CVprofile\else\def\CVprofile{sector-publico}\fi

% Debe compilarse con LuaLaTeX (la clase usa fontspec/fuentes OpenType)
\documentclass[localFont,alternative,10pt]{yaac-another-awesome-cv}

% -----------------------------------------------------------------------------
% Configuración compartida por todos los perfiles
% -----------------------------------------------------------------------------
% =============================================================================
% CONFIGURACIÓN: Información personal y de contacto
% Archivo: config/personal.tex
% Descripción: Fuente única de los datos personales del autor. Todos los
%              perfiles (profiles/*.tex) comparten esta configuración;
%              un cambio aquí se propaga a todas las variantes del CV.
% =============================================================================

% Nombre completo (usado en el pie de página y disponible para otras secciones)
\newcommand{\authorFullName}{Elmer Edison Achalma Mendoza}

% Definir el nombre del autor (Obligatorio)
% Uso: \name{<nombre>}{<apellido>}
\name{Elmer Edison}{Achalma Mendoza}

% Definir el lema del autor (Obligatorio)
% Uso: \tagline{<lema>}
\tagline{Economista | Docente | Análisis de Datos | Sector Público}

% Definir foto del autor (Opcional)
% Uso: \photo[<forma: circular, square, roundedsquare, squircle>]{<diámetro>}{<archivo>}
% La forma por defecto es circular
\photo{2.5cm}{../assets/images/profile/profile2}

% -----------------------------------------------------------------------------
% Redes sociales e información de contacto
% -----------------------------------------------------------------------------
\socialinfo{
	% LinkedIn (Opcional)
	\linkedin{achalmaedison}

	% GitHub (Opcional)
	\github{achalmed}\\

	% ORCID (Opcional) - Identificador único de investigador
	%\orcid{0000-0000-0000-0000}

	% Google Scholar (Opcional)
	%\googlescholar{<ID-google-scholar>}{Tu Nombre}

	% Sitio web personal (Opcional)
	\personalLink{achalmaedison.netlify.app}

	% Teléfono (Opcional)
	\smartphone{+51 934 179 301}\\

	% Email (Opcional)
	\email{elmer.achalma.09@unsch.edu.pe}

	% Email alternativo (Opcional)
	%\email{achalmaedison@gmail.com}\\

	% Dirección (Opcional)
	%\address{<dirección completa>}

	% Ubicación (Opcional)
	\infos{Ayacucho - Perú}

	% Redes sociales adicionales disponibles:
	% \viadeo{<usuario-viadeo>}
	% \stackoverflow{<id-stackoverflow>}
	% \stackexchange{<id-stackexchange>}
	% \twitter{<usuario-twitter>}
	% \skype{<usuario-skype>}
	% \researchgate{<usuario-researchgate>}
	% \medium{<usuario-medium>}
	% \kaggle{<usuario-kaggle>}
}

% -----------------------------------------------------------------------------
% Bibliografía (solo se imprime si el perfil activa la sección publicaciones)
% -----------------------------------------------------------------------------
\addbibresource{../bibliography/my_publications.bib}
   % datos personales, contacto, foto, bibliografía
% =============================================================================
% CONFIGURACIÓN: Ajustes de diseño y presentación
% Archivo: config/settings.tex
% Descripción: Ajustes visuales compartidos por todos los perfiles.
%              Las opciones de clase (localFont, alternative, tamaño de
%              fuente, compact) se definen en el \documentclass de cv.tex:
%
%   [localFont]      - Usar las fuentes incluidas en la carpeta fonts/
%   [alternative]    - Usar el diseño alternativo del encabezado
%   [10pt|11pt|12pt] - Tamaño de fuente (10pt por defecto)
%   [compact]        - Reducir el espacio vertical entre experiencias
% =============================================================================

% Cambiar el tamaño de la columna izquierda (por defecto 2.5cm)
% Uso: \setleftcolumnlength{<tamaño>}  — recomendado entre 1.5cm y 2.5cm
%\setleftcolumnlength{2.0cm}
   % ajustes de diseño

% Carga condicional de la sección de anexos: los perfiles la invocan con
% \inputAnexos{...}; no hace nada si \CVsinAnexos está definida (ANEXOS=0).
\newcommand{\inputAnexos}[1]{\ifdefined\CVsinAnexos\else\input{#1}\fi}

\begin{document}

	% Encabezado y pie de página
	\makecvheader
	\makecvfooter
	{\textsc{\today}}                        % Fecha (izquierda)
	{\textsc{\authorFullName{} --- CV}}      % Nombre (centro)
	{\thepage}                               % Número de página (derecha)

	% Cuerpo del CV: el perfil activo decide qué secciones se incluyen
	\input{profiles/\CVprofile}

\end{document}
"
%
% Si \CVprofile no está definida se usa el perfil por defecto:
\ifdefined\CVprofile\else\def\CVprofile{sector-publico}\fi

% Debe compilarse con LuaLaTeX (la clase usa fontspec/fuentes OpenType)
\documentclass[localFont,alternative,10pt]{yaac-another-awesome-cv}

% -----------------------------------------------------------------------------
% Configuración compartida por todos los perfiles
% -----------------------------------------------------------------------------
% =============================================================================
% CONFIGURACIÓN: Información personal y de contacto
% Archivo: config/personal.tex
% Descripción: Fuente única de los datos personales del autor. Todos los
%              perfiles (profiles/*.tex) comparten esta configuración;
%              un cambio aquí se propaga a todas las variantes del CV.
% =============================================================================

% Nombre completo (usado en el pie de página y disponible para otras secciones)
\newcommand{\authorFullName}{Elmer Edison Achalma Mendoza}

% Definir el nombre del autor (Obligatorio)
% Uso: \name{<nombre>}{<apellido>}
\name{Elmer Edison}{Achalma Mendoza}

% Definir el lema del autor (Obligatorio)
% Uso: \tagline{<lema>}
\tagline{Economista | Docente | Análisis de Datos | Sector Público}

% Definir foto del autor (Opcional)
% Uso: \photo[<forma: circular, square, roundedsquare, squircle>]{<diámetro>}{<archivo>}
% La forma por defecto es circular
\photo{2.5cm}{../assets/images/profile/profile2}

% -----------------------------------------------------------------------------
% Redes sociales e información de contacto
% -----------------------------------------------------------------------------
\socialinfo{
	% LinkedIn (Opcional)
	\linkedin{achalmaedison}

	% GitHub (Opcional)
	\github{achalmed}\\

	% ORCID (Opcional) - Identificador único de investigador
	%\orcid{0000-0000-0000-0000}

	% Google Scholar (Opcional)
	%\googlescholar{<ID-google-scholar>}{Tu Nombre}

	% Sitio web personal (Opcional)
	\personalLink{achalmaedison.netlify.app}

	% Teléfono (Opcional)
	\smartphone{+51 934 179 301}\\

	% Email (Opcional)
	\email{elmer.achalma.09@unsch.edu.pe}

	% Email alternativo (Opcional)
	%\email{achalmaedison@gmail.com}\\

	% Dirección (Opcional)
	%\address{<dirección completa>}

	% Ubicación (Opcional)
	\infos{Ayacucho - Perú}

	% Redes sociales adicionales disponibles:
	% \viadeo{<usuario-viadeo>}
	% \stackoverflow{<id-stackoverflow>}
	% \stackexchange{<id-stackexchange>}
	% \twitter{<usuario-twitter>}
	% \skype{<usuario-skype>}
	% \researchgate{<usuario-researchgate>}
	% \medium{<usuario-medium>}
	% \kaggle{<usuario-kaggle>}
}

% -----------------------------------------------------------------------------
% Bibliografía (solo se imprime si el perfil activa la sección publicaciones)
% -----------------------------------------------------------------------------
\addbibresource{../bibliography/my_publications.bib}
   % datos personales, contacto, foto, bibliografía
% =============================================================================
% CONFIGURACIÓN: Ajustes de diseño y presentación
% Archivo: config/settings.tex
% Descripción: Ajustes visuales compartidos por todos los perfiles.
%              Las opciones de clase (localFont, alternative, tamaño de
%              fuente, compact) se definen en el \documentclass de cv.tex:
%
%   [localFont]      - Usar las fuentes incluidas en la carpeta fonts/
%   [alternative]    - Usar el diseño alternativo del encabezado
%   [10pt|11pt|12pt] - Tamaño de fuente (10pt por defecto)
%   [compact]        - Reducir el espacio vertical entre experiencias
% =============================================================================

% Cambiar el tamaño de la columna izquierda (por defecto 2.5cm)
% Uso: \setleftcolumnlength{<tamaño>}  — recomendado entre 1.5cm y 2.5cm
%\setleftcolumnlength{2.0cm}
   % ajustes de diseño

% Carga condicional de la sección de anexos: los perfiles la invocan con
% \inputAnexos{...}; no hace nada si \CVsinAnexos está definida (ANEXOS=0).
\newcommand{\inputAnexos}[1]{\ifdefined\CVsinAnexos\else\input{#1}\fi}

\begin{document}

	% Encabezado y pie de página
	\makecvheader
	\makecvfooter
	{\textsc{\today}}                        % Fecha (izquierda)
	{\textsc{\authorFullName{} --- CV}}      % Nombre (centro)
	{\thepage}                               % Número de página (derecha)

	% Cuerpo del CV: el perfil activo decide qué secciones se incluyen
	\input{profiles/\CVprofile}

\end{document}
"
%
% Si \CVprofile no está definida se usa el perfil por defecto:
\ifdefined\CVprofile\else\def\CVprofile{sector-publico}\fi

% Debe compilarse con LuaLaTeX (la clase usa fontspec/fuentes OpenType)
\documentclass[localFont,alternative,10pt]{yaac-another-awesome-cv}

% -----------------------------------------------------------------------------
% Configuración compartida por todos los perfiles
% -----------------------------------------------------------------------------
% =============================================================================
% CONFIGURACIÓN: Información personal y de contacto
% Archivo: config/personal.tex
% Descripción: Fuente única de los datos personales del autor. Todos los
%              perfiles (profiles/*.tex) comparten esta configuración;
%              un cambio aquí se propaga a todas las variantes del CV.
% =============================================================================

% Nombre completo (usado en el pie de página y disponible para otras secciones)
\newcommand{\authorFullName}{Elmer Edison Achalma Mendoza}

% Definir el nombre del autor (Obligatorio)
% Uso: \name{<nombre>}{<apellido>}
\name{Elmer Edison}{Achalma Mendoza}

% Definir el lema del autor (Obligatorio)
% Uso: \tagline{<lema>}
\tagline{Economista | Docente | Análisis de Datos | Sector Público}

% Definir foto del autor (Opcional)
% Uso: \photo[<forma: circular, square, roundedsquare, squircle>]{<diámetro>}{<archivo>}
% La forma por defecto es circular
\photo{2.5cm}{../assets/images/profile/profile2}

% -----------------------------------------------------------------------------
% Redes sociales e información de contacto
% -----------------------------------------------------------------------------
\socialinfo{
	% LinkedIn (Opcional)
	\linkedin{achalmaedison}

	% GitHub (Opcional)
	\github{achalmed}\\

	% ORCID (Opcional) - Identificador único de investigador
	%\orcid{0000-0000-0000-0000}

	% Google Scholar (Opcional)
	%\googlescholar{<ID-google-scholar>}{Tu Nombre}

	% Sitio web personal (Opcional)
	\personalLink{achalmaedison.netlify.app}

	% Teléfono (Opcional)
	\smartphone{+51 934 179 301}\\

	% Email (Opcional)
	\email{elmer.achalma.09@unsch.edu.pe}

	% Email alternativo (Opcional)
	%\email{achalmaedison@gmail.com}\\

	% Dirección (Opcional)
	%\address{<dirección completa>}

	% Ubicación (Opcional)
	\infos{Ayacucho - Perú}

	% Redes sociales adicionales disponibles:
	% \viadeo{<usuario-viadeo>}
	% \stackoverflow{<id-stackoverflow>}
	% \stackexchange{<id-stackexchange>}
	% \twitter{<usuario-twitter>}
	% \skype{<usuario-skype>}
	% \researchgate{<usuario-researchgate>}
	% \medium{<usuario-medium>}
	% \kaggle{<usuario-kaggle>}
}

% -----------------------------------------------------------------------------
% Bibliografía (solo se imprime si el perfil activa la sección publicaciones)
% -----------------------------------------------------------------------------
\addbibresource{../bibliography/my_publications.bib}
   % datos personales, contacto, foto, bibliografía
% =============================================================================
% CONFIGURACIÓN: Ajustes de diseño y presentación
% Archivo: config/settings.tex
% Descripción: Ajustes visuales compartidos por todos los perfiles.
%              Las opciones de clase (localFont, alternative, tamaño de
%              fuente, compact) se definen en el \documentclass de cv.tex:
%
%   [localFont]      - Usar las fuentes incluidas en la carpeta fonts/
%   [alternative]    - Usar el diseño alternativo del encabezado
%   [10pt|11pt|12pt] - Tamaño de fuente (10pt por defecto)
%   [compact]        - Reducir el espacio vertical entre experiencias
% =============================================================================

% Cambiar el tamaño de la columna izquierda (por defecto 2.5cm)
% Uso: \setleftcolumnlength{<tamaño>}  — recomendado entre 1.5cm y 2.5cm
%\setleftcolumnlength{2.0cm}
   % ajustes de diseño

% Carga condicional de la sección de anexos: los perfiles la invocan con
% \inputAnexos{...}; no hace nada si \CVsinAnexos está definida (ANEXOS=0).
\newcommand{\inputAnexos}[1]{\ifdefined\CVsinAnexos\else\input{#1}\fi}

\begin{document}

	% Encabezado y pie de página
	\makecvheader
	\makecvfooter
	{\textsc{\today}}                        % Fecha (izquierda)
	{\textsc{\authorFullName{} --- CV}}      % Nombre (centro)
	{\thepage}                               % Número de página (derecha)

	% Cuerpo del CV: el perfil activo decide qué secciones se incluyen
	\input{profiles/\CVprofile}

\end{document}
